\scanpage{26}
\mspage{-\handwritten{3}}

Under these conditions one easily sees that $q.i=j.r$ (see for example Serre, algebraic groups and class fields, proof of lemma 11, page \handwritten{62} ; it is there that the hypothesis that the base scheme is the spectrum of a field comes in ; if one wants to generalize, one will therefore have to study \deleted{\unsure{h|m}}the morphism $q.i$).

One finally arrives at the commutative diagram
\diagram{
\begin{tikzcd}[ampersand replacement=\&,column sep=1.0em,row sep=2.5em]
\mathbb{A}
  \handwritten{\arrow[rr]}
  \arrow[rr,phantom,"\handwritten{\Delta}^{p}"]
  \handwritten{\arrow[dr,"a"]}
\& \& \handwritten{\otimes}^{p}\mathbb{A}
  \handwritten{\arrow[rr]}
  \arrow[rr,phantom,"m^p"]
  \handwritten{\arrow[dr,"q"']}
\& \& \mathbb{A}
\\
\& \handwritten{\Sigma}^{p}\mathbb{A}
  \handwritten{\arrow[ur,"i"']}
  \handwritten{\arrow[dr,"r"]}
\& \& S^p\mathbb{A}
  \handwritten{\arrow[ur,"b"]}
\&
\\
\& \& \mathbb{A}^{(p)}
  \handwritten{\arrow[ur,"j"]}
\& \&
\end{tikzcd}
}

Now $m^p.\handwritten{\Delta}^{p}$ is nothing other than the ring homomorphism corresponding to multiplication by $p$ in $G$. This homomorphism will be denoted by $\uline{p}$. Moreover, $b.j$ is nothing other than the Frobenius $F$. If we denote by $V$ and call Verschiebung the map $r.a$, we have $F.V=\uline{p}$. Finally, as the diagram is self-dual, $F$ and $V$ are interchanged in Cartier duality.

In the diagram above, $S^p\mathbb{A}$ is the only vector space not equipped with an algebra structure (it is a coalgebra, whereas $\handwritten{\Sigma}^{p}A$ is an algebra but not a coalgebra). If \unsure{nons|names|we} are not interested in duality and if we want to reason only in terms of \unsure{algebraic|algebraic} varieties, we therefore obtain the diagram
% archive sheet number: 25
